\addtocontents{toc}{\protect\setcounter{tocdepth}{0}}
\chapter[Anexo: Consideraciones Éticas]{ANEXO: Consideraciones Éticas}

\section{Evaluación ética del proyecto}

El proyecto se presenta al Comité de Ética Científica de la Universidad del
Bío-Bío mediante el protocolo abreviado para investigaciones de pregrado. Se
declara como una investigación aplicada de desarrollo tecnológico en nivel
TRL~3, que utiliza únicamente documentos existentes, como guías clínicas, el
programa PINDA y los contenidos entregados por el equipo de Enfermería. No
involucra participantes humanos, animales ni muestras biológicas.

La revisión del prototipo por parte de la contraparte de Enfermería forma parte
del trabajo colaborativo del proyecto y se limita a la pertinencia de los
contenidos y de los criterios de alerta. No implica aplicar instrumentos ni
recolectar datos personales. Toda evaluación futura con cuidadores queda a cargo
del equipo de Enfermería y requerirá su propia revisión ética.

\section{Alcance del sistema}

OnCuidar es una herramienta de apoyo educativo y de registro. No diagnostica, no
modifica tratamientos, no calcula dosis y no reemplaza al equipo de salud. Cada
alerta recuerda que la orientación es educativa y, ante un riesgo, recomienda
contactar al equipo tratante o acudir a urgencia.

\section{Protección de datos de salud}

Los datos de salud son datos sensibles según la Ley~19.628 sobre protección de
la vida privada. La aplicación los resguarda con cuatro mecanismos.

\begin{itemize}
    \item \emph{Cifrado en el dispositivo}: los datos personales del paciente,
    los registros clínicos, las conversaciones, los recordatorios y las listas
    se cifran con AES-256-GCM antes de salir del teléfono.
    \item \emph{Gestión de claves}: la clave de cada usuario se guarda en el
    almacenamiento seguro del teléfono y su copia en el servidor se sella con
    una clave maestra y secretos de entorno que no forman parte del
    repositorio. El servidor solo ve texto cifrado y no puede leer los datos de
    salud.
    \item \emph{Control de acceso}: cada solicitud exige una sesión firmada
    (JWT) y el servidor limita cada dato a su propietario, de modo que un
    recurso ajeno responde como inexistente, y rechaza cualquier campo sensible
    sin cifrar. Los enlaces de correo para verificar o recuperar la cuenta son
    de un solo uso y vencen en una hora.
    \item \emph{Minimización}: la aplicación solicita solo los datos necesarios
    para el seguimiento y no utiliza servicios de mensajería externos.
\end{itemize}

Los informes PDF y Excel se generan en el teléfono del cuidador y solo salen de
él si el propio cuidador decide compartirlos.

\section{Uso de herramientas de inteligencia artificial}

La aplicación no utiliza inteligencia artificial generativa para orientar al
cuidador. Las respuestas provienen de preguntas frecuentes validadas y la
clasificación de alertas de reglas explícitas. Cualquier incorporación futura de
este tipo de tecnología deberá someterse a una nueva revisión ética.
